\documentclass[10pt,a4paper]{amsart}
\usepackage[margin=19mm]{geometry}
\usepackage{amsmath,amssymb,mathtools}
\usepackage[scr=rsfs,cal=euler]{mathalfa}
\usepackage{xfrac}
\usepackage[unicode,hidelinks,bookmarks=false]{hyperref}
\newcommand{\ie}{i.e.\ }
\newcommand{\cf}{cf.\ }
\newcommand{\resp}{respectively\ }
\newcommand{\Subsection}{\subsection}
\providecommand{\nobreakdash}{\nobreak\hbox{-}}
\setlength{\emergencystretch}{2em}
\multlinegap0pt
\usepackage{bm}
\usepackage{bbm}
\usepackage{dsfont}
\usepackage{xspace}
\usepackage{cancel}
\usepackage{mathtools}
\usepackage{amsthm}
\makeatletter
\@ifundefined{lemma}{\newtheorem{lemma}{Lemma}}{}
\makeatother
\DeclareMathOperator*{\argmin}{argmin}
\DeclareMathOperator{\grad}{\nabla}
\newcommand{\abs}[1]{|#1|}
\newcommand{\norm}[2][]{{\|{#2}\|_{}}_{#1}}
\newcommand{\subjectto}{\text{\quad s.t.\quad }}
\newcommand{\udag}{u^{\dagger}}
\title[Denoising of image gradients and total generalized variation]{Denoising of image gradients and total generalized variation denoising}
\author{Birgit Komander, Dirk A. Lorenz, Lena Vestweber}
\date{Source: arXiv:1712.08585v3 (CC BY 4.0)}
\begin{document}
\maketitle

\noindent\emph{Source and licence.} Denoising of image gradients and total generalized variation denoising. Version arXiv:1712.08585v3 (CC BY 4.0), distributed under the Creative Commons Attribution 4.0 International licence, as recorded in the arXiv licence metadata for this exact version and shown on the abstract page https://arxiv.org/abs/1712.08585v3. Full rights notice: licenses/arxiv-1712.08585-CC-BY-4.0.txt.\par\smallskip

\noindent\textit{Editorial scope.} Counted unit: the Lemma whose source label is \texttt{lem:perfect-reconstruction-udag} in the article (source0.tex lines 189--199, source label \texttt{lem:perfect-reconstruction-udag}), delivered together with the article's own complete original proof of it (source0.tex lines 201--221, one \texttt{proof} environment of 21 lines). No mathematics is written, repaired or re-derived: statement and proof are the article's own text line for line. Required context retained uncounted: none. Every definition, display and label of those lines is retained unchanged. Omitted from this package and not redistributed: 1--188, 200--200, 222--1983. Nothing inside a retained excerpt is omitted or altered: every retained body is delivered exactly as the article's lines are written, so no omission receipt of any kind accompanies this package. Citations: the retained excerpts carry no citation command at all, so no bibliography entry of the article is transcribed and no reference is redistributed. Reference adaptation: none was needed and none was made; every reference inside the delivered unit resolves inside it, so no unresolved reference or citation is produced. Printed slips kept literal: no printed word, symbol, hypothesis or number of the counted statement or of its proof was altered. Layout and class: the article's own class is \texttt{svjour3} with packages including amsmath, amssymb, graphicx, subcaption, enumerate, times, hyperref, algorithm, algpseudocode, mathtools, float, pgfplots, booktabs, xcolor; the wrapper is the lane's pinned \texttt{amsart} editorial wrapper at 10pt on A4, which restates the theorem-like environments the retained text needs and adds the article's own macro definitions that the retained text needs (\texttt{\textbackslash abs}, \texttt{\textbackslash argmin}, \texttt{\textbackslash grad}, \texttt{\textbackslash norm}, \texttt{\textbackslash subjectto}, \texttt{\textbackslash udag}), with extra wrapper packages (bm, bbm, dsfont, xspace, cancel, mathtools). The delivered numbering is the unit's own. Assets: no retained span contains a figure environment, a graphics-inclusion command, a PDF-page inclusion, a table or a TikZ picture; no image file is copied into this package, the package declares no asset dependency and no image file is redistributed with it. Licence: version arXiv:1712.08585v3 of this article is distributed under the Creative Commons Attribution 4.0 International licence, as recorded in the arXiv licence metadata for this exact version and shown on the abstract page https://arxiv.org/abs/1712.08585v3. A full rights notice is delivered at licenses/arxiv-1712.08585-CC-BY-4.0.txt. Characters: every retained excerpt is pure ASCII, so the delivered engine, XeLaTeX with its default Latin Modern text font, reports no missing character.
\begin{lemma}
  \label{lem:perfect-reconstruction-udag}
  Assume that $\udag$ and $u_{0}$ fulfill $\int_{\Omega}\udag = \int_{\Omega}u_{0}$ and
  let $v = \grad \udag$ and $\delta_{1} = \norm[2]{\udag-u_{0}}$. Then it holds that
  \begin{equation}
    \label{eq:tvgrad-true-delta}
    \udag = \argmin_u \norm[1]{\abs{\grad u-v}}
    \subjectto\norm[2]{u-u_0}\leq \delta_{1},
  \end{equation}
  i.e. $\udag$ is the unique solution of the denoising problem.
\end{lemma}

\begin{proof}
  The set of minimizers is
  \begin{align*}
    \argmin_{u}\norm[1]{\abs{\grad u - \grad\udag}}
    \subjectto \norm[2]{u-u_0}
    \leq\norm[2]{\udag - u_0}.
  \end{align*}
  Clearly, $\udag$ is within this set, since the optimal value
  is $0$ and $\udag$ is feasible, because the constraint is 
  trivially fulfilled.
  
  To show that $\udag$ is indeed the unique solution, consider any other $u$ that also produces an objective value of zero.
  This implies $\grad u = \grad\udag$, i.e. $u = \udag + c$ for some constant $c$.
  Thus, $u$ fulfills the constraint $\norm[2]{u-u_{0}}^{2}\leq \norm[2]{\udag-u_{0}}^{2}$ if $\norm[2]{\udag-u_{0} + c}^{2}\leq \norm[2]{\udag-u_{0}}^{2}$.
  We expand the left hand side and get, writing $\abs{\Omega}$ for the measure of $\Omega$,
  \[
  \norm[2]{\udag-u_{0}}^{2} + 2c\int_{\Omega}(\udag-u_{0}) + c^{2}\abs{\Omega} \leq \norm[2]{\udag-u_{0}}^{2}.
  \]
  Since the middle integral vanishes by assumption, we see that $c=0$.\\
  \qed
\end{proof}

\end{document}
